%% Kernel fixture: a tensor standing where a curve meets another wire (the
%% corner cut of Section III, lines 468-497: a curve crosses two of a
%% tensor's indices and a small square stands on it at each meeting).
%% Expected records: three ket atoms with physical legs up; one string routed
%% north of the row with an over habit; two box atoms whose places are the
%% crossings of that string with the outer two legs.  The second picture is
%% the same production with two curves as its operands.
%% Expected ink: each box stands ON the meeting.  The curve keeps its whole
%% route through the box and the index keeps its whole run through it, so a
%% box is what sits where the two meet and not a junction of them: the middle
%% leg, which carries no box, shows the same meeting as a plain crossing.
\documentclass{standalone}
\usepackage{tenkz}
\makeatletter
\ExplSyntaxOn
\file_input:n { tenkz-kernel.code.tex }
\ExplSyntaxOff
\makeatother
\begin{document}
\tenkzkernel
\begin{tenkz}[rows={ket}, cols=3, physical=up]
  \tnwire[kind=string, species=g, name=g, crossing=over,
          route={n of (1,1) .. (1,3)}]{open}{open}
  \tn[skin=box, at=crossing of g and leg n of (1,1)]{}
  \tn[skin=box, at=crossing of g and leg n of (1,3)]{}
\end{tenkz}
\qquad
\begin{tenkz}[rows={wire,wire}, cols=4, bonds=none]
  \tnwire[kind=string, species=left,  name=a]{(1,1)}{(2,3)}
  \tnwire[kind=string, species=right, name=b,
          cross={under at crossing of a and b}]{(1,4)}{(2,2)}
  \tn[skin=box, at=crossing of a and b]{$T$}
\end{tenkz}
\end{document}
